\section{Experiments}

We ran six experiments (\cref{tab:experiments}). E1 measures the quality of the answers, E2 the
search alone, and E3--E6 the cloud side.

\begin{table}[ht]
\centering\small
\begin{tabularx}{\linewidth}{llX}
\toprule
 & Experiment & What we measure \\
\midrule
E1 & Accuracy & Systems A, A0, A2, C and N with Haiku 5.5, Qwen2.5-7B and Qwen2.5-3B, on our question
              set and on FinanceBench: correct, wrong and declined answers, tokens and cost \\
E2 & Retrieval & BM25, semantic and hybrid search without an LLM \\
E3 & Latency per step & Time of each step, with Qwen2.5-3B on the CPU of the cloud machine \\
E4 & Load test & Questions per second and latency for 1, 10 and 50 users, without the LLM \\
E5 & Spark scaling & Run time of both Spark jobs with 1 to 32 cores \\
E6 & Updates & Time to add newly filed 10-Qs, measured and estimated for the busiest days \\
\bottomrule
\end{tabularx}
\caption{Experiments.}
\label{tab:experiments}
\end{table}

\subsection{Question sets}

\paragraph{Our question set.} We wrote 60 questions about ten large companies: 3M, Coca-Cola,
Johnson \& Johnson, Caterpillar and Honeywell (2023 filings), and PepsiCo, Tesla, Abbott, IBM and
Lockheed Martin (2024 filings). 50 form the test set; the other 10 are development questions, used
only to set the search threshold. The test set has six categories (\cref{tab:question-examples}),
including questions that the system should decline. Every correct answer was verified against the
filing itself, not against our database.

\begin{table}[ht]
\centering\small
\begin{tabularx}{\linewidth}{lrX}
\toprule
Category & $n$ & Example (correct answer) \\
\midrule
lookup     & 10 & ``What was 3M's revenue in Q2 2023?'' (\$8,325 million) \\
change     & 10 & ``How much did 3M's net sales change year over year in the second quarter of 2023?''
                  ($-4.3\%$, from \$8,702 million) \\
paraphrase & 8  & ``How big was J\&J's top line in Q1 2023?'' (revenue in everyday words) \\
why        & 12 & ``Why did Coca-Cola's operating income increase in Q2 2023?'' (favourable pricing
                  initiatives and channel and package mix) \\
decline    & 8  & ``What was IKEA's revenue in Q2 2024?'' (not an SEC filer: must be declined) \\
latest     & 2  & ``What was Honeywell's EPS?'' (no period: the latest quarter, with the assumption stated) \\
\bottomrule
\end{tabularx}
\caption{The six categories of our test set, with one example each.}
\label{tab:question-examples}
\end{table}

\paragraph{FinanceBench.} The 150 public FinanceBench questions~\cite{financebench}, written by
analysts about 84 filings, mostly annual reports. They come in three types of 50: calculations (``What
is Amazon's FY2017 days payable outstanding?''), questions that need an analyst's judgement (``Is 3M a
capital-intensive business based on FY2022 data?'') and other questions. As in the benchmark's
standard setting, every system is given the correct filing.

\subsection{Systems and models}

We run the workflow (A), its variants A0 (template answer) and A2 (with the financial statements),
simple RAG (C) and closed book (N). The hosted model is Claude Haiku 5.5. The local models are
Qwen2.5-7B and Qwen2.5-3B, run on a MacBook and connected to the cloud database; for E3, Qwen2.5-3B
also runs on the CPU of the cloud machine. To see how much a much larger model knows without any
documents, we also run N with Claude Opus 5.5. E3--E6 ran on the c7i.8xlarge. For E4 the LLM steps
are replaced by fixed outputs, so that only our own code is measured.

\subsection{Grading}
\label{sec:grading}

\paragraph{Numbers} are graded by code. An answer is correct if it states the correct value to the
precision it gives: ``\$8.3 billion'' is correct for \$8,325 million, ``\$8.4 billion'' is not. A
change must be within 0.15 percentage points.

\paragraph{Free text} (``why'' answers and all FinanceBench answers) is graded by an LLM judge
(Claude Haiku 5.5). It sees the question, the correct answer and the system's answer and decides
whether the answer is correct, incorrect or declined. To check the judge, a stronger model (Claude
Sonnet 5.5) graded a random sample of 300 of these answers again (\cref{tab:grader}). The two agree on
89\% of the answers, and on 95\% for the distinction that matters for accuracy, correct or not. Cohen's
$\kappa$, which corrects agreement for agreement by chance (1 is perfect, 0 is chance), is 0.84,
usually read as strong agreement.

\begin{table}[ht]
\centering\small
\begin{tabular}{lrrr}
\toprule
Judge (Haiku) $\downarrow$ \quad Sonnet $\rightarrow$ & correct & incorrect & declined \\
\midrule
correct   & \textbf{125} & 8  & 1 \\
incorrect & 4   & \textbf{74} & 18 \\
declined  & 1   & 0  & \textbf{69} \\
\bottomrule
\end{tabular}
\caption{Grader check on 300 random LLM-graded answers: verdicts of our judge (rows) against a
stronger model (columns).}
\label{tab:grader}
\end{table}

\subsection{Metrics}
\label{sec:metrics}

\begin{itemize}
  \item \textbf{Accuracy:} the share of questions answered correctly. A refusal counts as a decline:
        it is not correct, but it is reported separately from a \textbf{wrong answer}, because a
        wrong number is worse than none. \textbf{Wrong declines} are refusals of questions that could
        have been answered.
  \item \textbf{Repeated runs and confidence intervals.} Haiku gives slightly different answers each
        time, so every Haiku result was run three times and we report the mean $\pm$ the standard
        deviation. Qwen always gives the same answer, so it ran once. Because the test sets are
        small, we also report a 95\% confidence interval, computed by bootstrapping: we draw the
        questions with replacement 2,000 times and take the range that contains 95\% of the
        resulting accuracies.
  \item \textbf{Tokens and cost:} the number of text pieces (tokens) sent to and produced by the LLM per
        question, and the API cost per 1,000 questions.
  \item \textbf{Retrieval (E2):} \emph{hit@k} is the share of questions for which at least one correct
        passage is among the first $k$ search results. \emph{MRR} (mean reciprocal rank) averages
        $1/r$, where $r$ is the position of the first correct passage: 1 if it is first, 0.5 if
        second, 0 if it is not in the top 10.
  \item \textbf{Latency and throughput (E3, E4):} the median time per question (half of the questions
        are faster), the 95th percentile (95\% are faster), and the number of questions answered per
        second.
  \item \textbf{Speed-up (E5):} the run time on one core divided by the run time on $n$ cores; ideal
        scaling gives a speed-up of $n$.
\end{itemize}
